\section{Tecnologías de baterías}\label{sec:tecnologias}

Una batería recargable almacena energía mediante una reacción de reducción-oxidación reversible. Durante la descarga, el electrodo negativo se oxida y libera electrones que circulan por el circuito externo hasta el electrodo positivo, que se reduce; el electrolito permite el paso de iones entre ambos electrodos. Durante la carga, una fuente externa invierte la reacción \cite{linden}. El proceso no es completamente reversible: en cada ciclo parte del material activo deja de participar en la reacción, y por eso la capacidad disminuye con el uso.

Entre las tecnologías de almacenamiento electroquímico estudiadas se encuentran las baterías de plomo-ácido, ion-litio, níquel-cadmio, sodio-azufre y de flujo. Las tres últimas se descartan para el proyecto. Las de níquel-cadmio usan cadmio, un metal tóxico, y su costo instalado es de 1200 a \SI{1500}{USD/kWh}, el doble o más que el del plomo-ácido \cite{may2018}. Las de sodio-azufre operan a más de \SI{300}{\celsius} y las de flujo almacenan el electrolito en tanques externos; ambas están orientadas a instalaciones de gran escala de la red eléctrica \cite{may2018}, no a un banco residencial de \SI{24}{V}. El análisis se concentra entonces en el plomo-ácido, con sus tres variantes, y en el ion-litio de fosfato de hierro.

\subsection{Plomo-ácido}

La batería de plomo-ácido usa dióxido de plomo (PbO\textsubscript{2}) en el electrodo positivo, plomo esponjoso (Pb) en el negativo y una solución de ácido sulfúrico como electrolito. Durante la descarga ambos electrodos se convierten en sulfato de plomo y el ácido se consume, según
\begin{equation}
\mathrm{Pb + PbO_2 + 2H_2SO_4 \rightleftharpoons 2PbSO_4 + 2H_2O}.
\end{equation}
La tensión nominal de cada celda es de \SI{2}{V}; una batería de \SI{12}{V} tiene seis celdas en serie y un banco de \SI{24}{V} necesita doce \cite{linden}, \cite{pavlov}.

Es la tecnología recargable más madura. Se usa en vehículos, sistemas de respaldo y telecomunicaciones, su costo de batería es de 200 a \SI{500}{USD/kWh} y se recicla casi en su totalidad \cite{may2018}. Sus desventajas son el peso y la vida en ciclos profundos. Su densidad de energía es de 35 a \SI{40}{Wh/kg}, frente a 150 a \SI{180}{Wh/kg} del ion-litio \cite{may2018}. Si la batería permanece descargada, el sulfato de plomo forma cristales grandes que ya no se reconvierten durante la carga; este proceso, llamado sulfatación, reduce la capacidad de forma permanente \cite{pavlov}. Por esa razón la descarga se limita: en el curso se usa una DOD máxima del 70\,\%, y muchos fabricantes declaran sus ciclos al 50\,\%. La práctica recomendada IEEE~1013 trata el dimensionamiento de este tipo de baterías en sistemas fotovoltaicos aislados \cite{ieee1013}.

Las variantes de plomo-ácido se diferencian en cómo se contiene el electrolito.

\subsubsection{Electrolito inundado}

En la batería inundada el electrolito es líquido y cubre libremente las placas. Al final de la carga parte de la corriente descompone el agua en hidrógeno y oxígeno, que salen por los tapones de ventilación. Esto obliga a reponer agua destilada de forma periódica, a instalar la batería en posición vertical y a ventilar el lugar para que el hidrógeno no se acumule \cite{linden}. Además, el ácido tiende a concentrarse en la parte inferior de la celda, por lo que el regulador debe aplicar cargas de ecualización: una sobrecarga controlada que produce gas y mezcla el electrolito \cite{pavlov}.

A cambio de ese mantenimiento, las baterías inundadas de ciclo profundo, como las de placas gruesas para carritos de golf o las de placa tubular OPzS, tienen el menor costo por Ah de la familia y toleran bien las sobrecargas, porque el agua perdida se puede reponer \cite{pavlov}. Para una vivienda, el mantenimiento y la necesidad de un espacio ventilado son las principales limitaciones.

\subsubsection{Regulada por válvula con separador AGM}

Las baterías reguladas por válvula (VRLA) son selladas. El oxígeno que se produce en el electrodo positivo al final de la carga se difunde hasta el negativo y se recombina para formar agua otra vez, de modo que la pérdida de agua es mínima y no hay que reponerla. Una válvula se abre solo si la presión interna supera un límite \cite{linden}, \cite{pavlov}.

En la variante AGM (\emph{absorbent glass mat}), el electrolito está absorbido en un separador de microfibra de vidrio colocado entre las placas. El separador no está saturado del todo, y los poros libres son los que permiten el paso del oxígeno hacia el electrodo negativo \cite{pavlov}. Como el electrolito queda retenido y muy cerca de las placas, la resistencia interna es baja y la batería entrega corrientes altas. No requiere reposición de agua y se puede instalar en otras posiciones, pero es sensible a la sobrecarga y a la temperatura alta: si la tensión de carga es excesiva, el gas que se libera por la válvula ya no se recupera y la batería se seca \cite{linden}. Por eso el regulador debe configurarse con las tensiones de carga que indica el fabricante.

\subsubsection{Regulada por válvula con electrolito en gel}

En la batería de gel, el ácido sulfúrico se mezcla con sílice pirogénica y forma un gel que inmoviliza el electrolito. Con el tiempo el gel se agrieta, y esas grietas son las que permiten la recombinación del oxígeno \cite{pavlov}. Igual que la AGM, es sellada y no requiere reposición de agua.

El gel evita que el ácido se estratifique y protege mejor a las placas en descargas profundas, por lo que estas baterías resisten más ciclos que las AGM en aplicaciones de ciclo diario, como las solares \cite{pavlov}. A cambio, su resistencia interna es mayor, entregan menos corriente pico y admiten una corriente de carga menor. Las versiones de placa tubular, conocidas como OPzV, son las de mayor vida en ciclos dentro del plomo-ácido sellado \cite{pavlov}.

\subsection{Ion-litio de fosfato de hierro (LiFePO\textsubscript{4})}

En una batería de ion-litio, los iones de litio se desplazan entre dos electrodos que los alojan dentro de su estructura cristalina: durante la carga salen del electrodo positivo y se insertan en el negativo, de grafito, y durante la descarga hacen el recorrido inverso \cite{tarascon2001}. El electrolito es una sal de litio disuelta en un solvente orgánico. Como el litio no se deposita como metal en ningún momento del ciclo normal, la reacción es muy reversible y la eficiencia de ida y vuelta es alta.

La química del electrodo positivo define el tipo de batería. Las de óxido de cobalto o de níquel-manganeso-cobalto (NMC) tienen mayor densidad de energía y se usan en electrónica portátil y vehículos eléctricos. La de fosfato de hierro y litio, LiFePO\textsubscript{4} o LFP, fue propuesta por Padhi \emph{et al.}\ en 1997 \cite{padhi1997}. En su estructura de olivino, el oxígeno está unido con fuerza al fósforo, lo que hace más difícil que se libere al calentarse; por eso las celdas LFP son más estables térmicamente que las de óxidos y tienen menor riesgo de incendio \cite{zubi2018}. Además, no usan cobalto ni níquel. La tensión nominal de cada celda es de \SI{3.2}{V}, de modo que ocho celdas en serie forman una batería de \SI{25.6}{V}, que es la versión comercial de \SI{24}{V} \cite{linden}.

Para el uso en sistemas aislados, las ventajas de la LFP son la profundidad de descarga y la vida en ciclos. Los fabricantes declaran ciclos con descargas del 80 al 100\,\%, frente al 50\,\% habitual en el plomo-ácido, y su densidad de energía es del orden de cuatro veces la del plomo-ácido \cite{may2018}, \cite{zubi2018}. No requiere mantenimiento ni ecualización, y su autodescarga es menor que la del plomo-ácido \cite{linden}.

Sus limitaciones son tres. Primero, cada batería necesita un sistema de gestión (BMS, \emph{battery management system}) que equilibre la tensión de las celdas y la desconecte por sobretensión, subtensión, sobrecorriente o temperatura fuera de rango \cite{linden}. Segundo, no se debe cargar por debajo de \SI{0}{\celsius}, porque el litio se deposita como metal sobre el grafito y daña la celda \cite{linden}. Tercero, la curva de descarga es muy plana, por lo que el estado de carga no se puede estimar solo con la tensión y se requiere un monitor que cuente la carga entrante y saliente. En costo, May \emph{et al.} reportan de 340 a \SI{590}{USD/kWh} para baterías LFP de la red eléctrica \cite{may2018}; los precios de baterías comerciales de \SI{24}{V} se revisan en la Sección~\ref{sec:opciones}.

\subsection{Síntesis}

La Tabla~\ref{tab:tecnologias} resume las diferencias cualitativas entre las cuatro tecnologías. Los valores numéricos de cada modelo comercial, tomados de sus fichas técnicas, se comparan en la Sección~\ref{sec:matriz}. Para todas ellas, la norma IEC~61427-1 define los ensayos de baterías destinadas a sistemas fotovoltaicos aislados \cite{iec61427}.

\begin{table}[!t]
\caption{Síntesis de las tecnologías consideradas}
\label{tab:tecnologias}
\centering
\footnotesize
\setlength{\tabcolsep}{2.5pt}
\begin{tabular}{>{\raggedright\arraybackslash}p{0.21\columnwidth} >{\raggedright\arraybackslash}p{0.17\columnwidth} >{\raggedright\arraybackslash}p{0.17\columnwidth} >{\raggedright\arraybackslash}p{0.17\columnwidth} >{\raggedright\arraybackslash}p{0.18\columnwidth}}
\toprule
 & Inundada & AGM & Gel & LiFePO\textsubscript{4} \\
\midrule
Tensión de celda & \SI{2}{V} & \SI{2}{V} & \SI{2}{V} & \SI{3.2}{V} \\
Celdas para \SI{24}{V} & 12 & 12 & 12 & 8 \\
Electrolito & Líquido libre & Absorbido en fibra de vidrio & Gelificado con sílice & Sal de litio en solvente orgánico \\
Reposición de agua & Sí & No & No & No \\
Ecualización & Sí & No & No & No \\
Ventilación & Necesaria & Reducida & Reducida & No \\
Requiere BMS & No & No & No & Sí \\
Principal limitación & Manteni\-miento & Sensible a sobrecarga & Corriente pico menor & Costo inicial; no cargar bajo \SI{0}{\celsius} \\
\bottomrule
\end{tabular}
\end{table}
